\subsection*{A correlated-radical separation lemma}
Let $F,D,Q,G,b,c\in\mathbb Z$, with $F\equiv D\equiv1\pmod3$ and
$Q\equiv0\pmod3$, and put $\delta=Q-1$. Suppose algebraic numbers satisfy
\[
 f^2=F,\qquad \tau^2=D,\qquad g^2=G,\qquad
 \delta t^2=Q(fb-c)^2-1.
\]
If $\rho=f+3\tau+9t+27g\in\mathbb Q$, then all four roots are rational.
Thus $f,\tau,g$ are integers; if $Q=0$ or $Q>1$, then $t$ is also an
integer. The signs are unrestricted, and the dependence of $t$ on $f$
is intentional. No independence of square classes is assumed.

\paragraph{Proof.}
Let $R=\mathbb Z_{(3)}$ and let $S$ be its integral closure in a finite
normal extension of $\mathbb Q$ containing all four roots. The ring $S$
is invariant under conjugation and $S\cap\mathbb Q=R$. Indeed, a reduced
rational with denominator divisible by $3$ cannot satisfy a monic equation
with coefficients whose denominators are prime to $3$: clearing denominators
would contradict reduction modulo $3$. The elements $f,\tau,g$ are integral
over $R$. Since $\delta$ is a unit of $R$, the radicand of $t$ belongs to
$S$; transitivity of integrality then gives $t\in S$.

If $f$ is irrational, an automorphism $\sigma$ of the normal extension
sends $f$ to $-f$. Subtracting the conjugate rational relation gives
\[
 2f=3(\sigma\tau-\tau)+9(\sigma t-t)+27(\sigma g-g).
\]
Consequently $f/3\in S$, since $2$ is a unit in $R$. Squaring gives
$F/9\in S\cap\mathbb Q=R$, contradicting $F\equiv1\pmod3$.
Thus $f\in\mathbb Q$, hence $f\in\mathbb Z$. Apply the same argument to
$(\rho-f)/3=\tau+3t+9g$ to obtain $\tau\in\mathbb Z$.
Now $t^2=[Q(fb-c)^2-1]/\delta\in R$ is rational and congruent to $1$
modulo $3$. Applying the argument once more to
$(\rho-f-3\tau)/9=t+3g$ gives $t\in\mathbb Q$, and then
$g\in\mathbb Q$, hence $g\in\mathbb Z$.

For the final integrality assertion, $Q=0$ gives $t^2=1$. If $Q>1$, set
$u=fb-c\in\mathbb Z$ and write $|t|=a/m$ in lowest terms. The equation
$Qu^2-(Q-1)t^2=1$ gives $m^2\mid Q-1$. The cases $u=0$ and $|u|=1$
are respectively impossible and integral. Otherwise reflect signs so
$u>1,t>0$; then $u<t<Qu/(Q-1)$. The norm-preserving transformation
\[
 u'=(2Q-1)u-2(Q-1)t,\qquad t'=(2Q-1)t-2Qu
\]
has $u'\in\mathbb Z$, $0<|u'|<u$, and leaves the reduced denominator
$m$ unchanged, since $2Q-1\equiv1\pmod m$. Descent to $|u|=1$ proves
$m=1$.\hfill$\square$

\paragraph{Scope.}
This is an ordinary algebraic lemma, including repeated or negative roots.
It is not a new end-to-end consistency calibration. An elimination product
must use the two $t$ roots belonging to \emph{each} selected sign of $f$;
a four-independent-radical formula with a fixed $t$ radicand would be wrong.
